\begin{textsample}{POS Dim 2 – mg – Score 26.00 – mg/mg\_em\_105.txt}  \label{ex:f2_pos_194}
5.3 Princípios e fundamentos da Educação Profissional e Técnica como Itinerário \textbf{formativo} O termo “ \textbf{itinerário} ” nos remete a caminho, roteiro, percurso.
Os \textbf{itinerários} \textbf{formativos} são possibilidades de caminhos para formação básica dos jovens e devem estar \textbf{alinhados} aos seus anseios, desejos, contextos e necessidades.
Nesse sentido, a \textbf{oferta} de um \textbf{itinerário} \textbf{formativo} estruturado pela Educação Profissional – reconhecido como o 5º \textbf{itinerário} – se apresenta como mais uma possibilidade de percurso \textbf{formativo} que pode ser trilhado no Ensino Médio e que busca oferecer estímulos e condições para que os jovens possam estruturar, consolidar e aplicar saberes e fazeres relacionados ao mundo do trabalho em todas as suas interfaces sociais e culturais.
A Resolução CNE/CP nº 1, de 05 de janeiro de 2021, que Define as \textbf{Diretrizes} Curriculares Nacionais Gerais para a Educação Profissional e Tecnológica, orienta em seu \textbf{capítulo} 3, art. 5º : Os \textbf{cursos} de Educação Profissional e Tecnológica podem ser organizados por \textbf{itinerários} \textbf{formativos}, observadas as orientações oriundas dos eixos tecnológicos. § 1º Os eixos tecnológicos deverão observar as distintas segmentações tecnológicas abrangidas, de forma a promover orientações específicas que sejam capazes de orientar as tecnologias contempladas em cada uma das distintas áreas tecnológicas identificadas. § 2º A não identificação de distintas áreas tecnológicas preservará as mesmas orientações dos eixos tecnológicos. § 3º O \textbf{Catálogo} Nacional de \textbf{Cursos} \textbf{Técnicos} ( CNCT ) e o \textbf{Catálogo} Nacional de \textbf{Cursos} Superiores de Tecnologia ( CNCST ) orientam a organização dos \textbf{cursos} dando visibilidade às \textbf{ofertas} de Educação Profissional e Tecnológica. § 4º O \textbf{itinerário} \textbf{formativo} deve contemplar a articulação de \textbf{cursos} e programas, configurando \textbf{trajetória} educacional consistente e programada, a partir de : I - estudos sobre os \textbf{itinerários} de profissionalização praticados no mundo do trabalho ; II - estrutura sócio-ocupacional da área de atuação \textbf{profissional} ; e III - fundamentos científico-tecnológicos dos processos produtivos de bens ou serviços. § 5º Entende -se por \textbf{itinerário} \textbf{formativo} na Educação Profissional e Tecnológica o conjunto de unidades curriculares, etapas ou módulos que compõem a sua organização em eixos tecnológicos e respectiva área tecnológica, podendo ser : I - propiciado internamente em um mesmo \textbf{curso}, mediante sucessão de unidades curriculares, etapas ou módulos com terminalidade ocupacional ; II - propiciado pela instituição educacional, mas construído horizontalmente pelo estudante, mediante unidades curriculares, etapas ou módulos de \textbf{cursos} diferentes de um mesmo eixo tecnológico e respectiva área tecnológica ; e III - construído verticalmente pelo estudante, propiciado ou não por instituição educacional, mediante sucessão progressiva de \textbf{cursos} ou \textbf{certificações} obtidas por avaliação e por reconhecimento de competências, desde a formação inicial até a pós-graduação tecnológica. § 6º Os \textbf{itinerários} \textbf{formativos} \textbf{profissionais} devem possibilitar um contínuo e articulado aproveitamento de estudos e de experiências \textbf{profissionais} devidamente avaliadas, reconhecidas e certificadas por instituições e \textbf{redes} de Educação Profissional e Tecnológica, criadas nos termos da \textbf{legislação} \textbf{vigente}. § 7º Os \textbf{itinerários} \textbf{formativos} \textbf{profissionais} podem ocorrer dentro de um \textbf{curso}, de uma área tecnológica ou de um eixo tecnológico, de modo a favorecer a verticalização da formação na Educação Profissional e Tecnológica, possibilitando, quando possível, diferentes percursos \textbf{formativos}, incluindo programas de aprendizagem \textbf{profissional}, observada a \textbf{legislação} trabalhista pertinente. § 8º Entende -se por eixo tecnológico a estrutura de organização da Educação Profissional e Tecnológica, considerando as diferentes matrizes tecnológicas nele existentes, por meio das quais são promovidos os agrupamentos de \textbf{cursos}, levando em consideração os fundamentos científicos que as sustentam, de forma a orientar o Projeto Pedagógico do \textbf{Curso} ( PPC ), identificando o conjunto de conhecimentos, habilidades, atitudes, valores e emoções que devem orientar e integrar a organização curricular, dando identidade aos respectivos \textbf{perfis} \textbf{profissionais}. ( Resolução CNE/CP nº 1, de 05 de janeiro de 2021, Art. 5º.)[27 ] Dessa forma, a \textbf{oferta} da formação \textbf{profissional} e tecnológica, em Minas Gerais, como possibilidade de \textbf{itinerário} \textbf{formativo} para estudantes do Ensino Médio, está sustentada em fundamentos \textbf{ético-políticos}, epistemológicos e didático-pedagógicos norteadores das práticas educativas voltadas à estruturação de uma formação humana e acadêmica de \textbf{qualidade}, pautada na inovação pedagógica, na valorização das aprendizagens experienciais e na superação da dicotomia entre a teoria e a prática.
As possibilidades de composição do 5º \textbf{itinerário} estão estruturadas em princípios norteadores que apontam para a \textbf{garantia} da formação integral dos estudantes, incluindo nela a construção e a compreensão dos conhecimentos historicamente escolarizados e organizados nas áreas do conhecimento e nos componentes curriculares da BNCC, o desenvolvimento de conhecimentos, de habilidades e de competências necessários para atuação no mundo do trabalho e em ocupações \textbf{técnicas} específicas e, também, à estruturação de valores e atitudes pautados na solidariedade, na criatividade, na inovação, na empatia e na colaboração. [ ^27 ] : BRASIL.
Ministério da Educação.
Conselho Nacional da Educação.
Câmara da Educação Básica.
Resolução CNE/CP nº 01, de 05 de janeiro de 2021.
Define as \textbf{Diretrizes} Curriculares Nacionais Gerais para a Educação Profissional e Tecnológica.
Disponível em : https://www.in.gov.br/web/dou/-/resolucao-cne/cp-n-1-de-5-de-janeiro-de-2021-297767578.
Acesso em 28 de março de 2021.
Neste sentido, reafirmamos os princípios norteadores da Educação Profissional e Tecnológica presentes nas \textbf{Diretrizes} Curriculares Nacionais Gerais para a Educação Profissional e Tecnológica ( DCNEPT ) : I - articulação com o setor produtivo para a construção coerente de \textbf{itinerários} \textbf{formativos}, com vista ao preparo para o exercício das profissões operacionais, \textbf{técnicas} e tecnológicas, na perspectiva da inserção laboral dos estudantes ; II - respeito ao princípio constitucional do pluralismo de ideias e de concepções pedagógicas ; III - respeito aos valores estéticos, políticos e éticos da educação nacional, na perspectiva do pleno desenvolvimento da pessoa, seu preparo para o exercício da cidadania e sua \textbf{qualificação} para o trabalho ; IV - centralidade do trabalho assumido como princípio educativo e base para a organização curricular, visando à construção de competências \textbf{profissionais}, em seus objetivos, conteúdos e estratégias de ensino e aprendizagem, na perspectiva de sua integração com a ciência, a cultura e a tecnologia ; V - estímulo à adoção da pesquisa como princípio pedagógico presente em um processo \textbf{formativo} voltado para um mundo permanentemente em transformação, integrando saberes cognitivos e socioemocionais, tanto para a produção do conhecimento, da cultura e da tecnologia, quanto para o desenvolvimento do trabalho e da intervenção que promova impacto social ; VI - a tecnologia, enquanto expressão das distintas formas de aplicação das bases científicas, como fio condutor dos saberes essenciais para o desempenho de diferentes funções no setor produtivo ; VII - indissociabilidade entre educação e prática social, bem como entre saberes e fazeres no processo de ensino e aprendizagem, considerando -se a historicidade do conhecimento, valorizando os sujeitos do processo e as metodologias ativas e inovadoras de aprendizagem centradas nos estudantes ; VIII - interdisciplinaridade assegurada no planejamento curricular e na prática pedagógica, visando à superação da \textbf{fragmentação} de conhecimentos e da segmentação e descontextualização curricular ; IX - utilização de estratégias educacionais que permitam a contextualização, a flexibilização e a interdisciplinaridade, favoráveis à compreensão de significados, garantindo a indissociabilidade entre a teoria e a prática \textbf{profissional} em todo o processo de ensino e aprendizagem ; X - articulação com o desenvolvimento socioeconômico e os arranjos produtivos locais ; XI - \textbf{observância} às necessidades específicas das pessoas com deficiência, Transtorno do Espectro Autista ( TEA ) e altas habilidades ou superdotação, gerando oportunidade de participação plena e efetiva em igualdade de condições no processo educacional e na sociedade ; XII - \textbf{observância} da condição das pessoas em \textbf{regime} de acolhimento ou internação e em \textbf{regime} de privação de liberdade, de maneira que possam ter acesso às \textbf{ofertas} educacionais, para o desenvolvimento de competências \textbf{profissionais} para o trabalho ; XIII - reconhecimento das identidades de gênero e étnico-raciais, assim como dos povos indígenas, quilombolas, populações do campo, imigrantes e itinerantes ; XIV - reconhecimento das diferentes formas de produção, dos processos de trabalho e das culturas a elas subjacentes, requerendo formas de ação diferenciadas ; XV - autonomia e flexibilidade na construção de \textbf{itinerários} \textbf{formativos} \textbf{profissionais} diversificados e \textbf{atualizados}, segundo interesses dos sujeitos, a relevância para o contexto local e as possibilidades de \textbf{oferta} das instituições e \textbf{redes} que oferecem Educação Profissional e Tecnológica, em consonância com seus respectivos projetos pedagógicos ; XVI - identidade dos \textbf{perfis} \textbf{profissionais} de conclusão de \textbf{curso}, que contemplem as competências \textbf{profissionais} requeridas pela natureza do trabalho, pelo desenvolvimento tecnológico e pelas demandas sociais, econômicas e ambientais ; XVII - autonomia da instituição educacional na concepção, elaboração, execução, avaliação e revisão do seu Projeto Político Pedagógico ( PPP ), construído como instrumento de referência de trabalho da comunidade escolar, respeitadas a \textbf{legislação} e as normas educacionais, estas \textbf{Diretrizes} Curriculares Nacionais e as \textbf{Diretrizes} complementares de cada sistema de ensino ; XVIII - \textbf{fortalecimento} das estratégias de colaboração entre os \textbf{ofertantes} de Educação Profissional e Tecnológica, visando ao maior alcance e à efetividade dos processos de ensino-aprendizagem, contribuindo para a empregabilidade dos egressos ; e XIX - promoção da inovação em todas as suas vertentes, especialmente a tecnológica, a social e a de processos, de maneira incremental e operativa. ( Resolução CNE/CP nº 1, de 05 de janeiro de 2021, Art. 3º ) No desenvolvimento das ações e estratégias pedagógicas, é imprescindível considerar o impacto das novas tecnologias e das novas formas de produção e transmissão das informações e dos conhecimentos, que demandam que os professores sejam mediadores, facilitadores dos processos de construção e estruturação de conhecimentos significativos, estimulando os estudantes a realizarem pesquisas e a agir de forma inovadora e protagonista em todos os espaços sociais.
As metodologias devem estar comprometidas com o desenvolvimento do espírito científico, com a integração dos conhecimentos e com a formação do sujeito-cidadão conectado com a realidade do mundo do trabalho.
Nesse sentido, as atividades precisam estar contextualizadas e contemplar experiências relacionadas ao mundo do trabalho e à prática \textbf{profissional} intrínseca à formação \textbf{técnica}.
Para além do conhecimento \textbf{técnico} e da utilização de equipamentos e materiais, é preciso propiciar aos estudantes condições para que, ao longo da vida, saibam interpretar, analisar, criticar, refletir, buscar soluções e propor alternativas potencializadas pela investigação e pelo reconhecimento dos territórios e dos contextos sociais.
Desse modo, serão formados sujeitos capazes de identificar demandas e problemas individuais, locais e globais.
E de mobilizar saberes, habilidades e competências para a estruturação de soluções baseadas em princípios éticos, justos, sustentáveis, solidários e colaborativos.

% matched lemmas: alinhar, atualizar, capítulo, catálogo, certificação, curso, diretriz, formativo, fortalecimento, fragmentação, garantia, itinerário, legislação, observância, oferta, ofertante, perfil, profissional, qualidade, qualificação, rede, regime, trajetória, técnico, vigente, ético-político
\end{textsample}
